\documentclass{beamer}

\usetheme{Madrid}
\usepackage{graphics}
\usepackage{graphicx}
\usepackage{epsfig}
\usepackage{pdfpages}


\usepackage{graphicx} % Allows including images
\usepackage{booktabs} % Allows the use of \toprule, 

% Custom font sizes for better readability
\setbeamerfont{title}{size=\Large}
\setbeamerfont{subtitle}{size=\large}
\setbeamerfont{author}{size=\normalsize}
\setbeamerfont{institute}{size=\normalsize}
\setbeamerfont{date}{size=\small}
\setbeamerfont{frametitle}{size=\Large}
\setbeamertemplate{footline}{}

\title{Startup Success Predictor using Multi Agent System}
\author[Group 8]{
\textbf{Group 8}\\[6pt]
\begin{tabular}{ll}
Darwin Thomas & 22CSB13 MDL22CS065 \\
K V Varun Krishnan & 22CSB30 MDL22CS115 \\
Manu S Panicker & 22CSB31 MDL22CS119 \\
M Arjun Krishna & 22CSB32 MDL22CS122
\end{tabular}
}




\institute{Model Engineering College}
\date{4 August 2025}
% Custom footer for all slides
\setbeamertemplate{footline}{
  \leavevmode%
  \hbox{%
    \begin{beamercolorbox}[wd=.8\paperwidth,ht=2.5ex,dp=1ex,left]{author in head/foot}%
      \hspace*{1em}\insertshortauthor
    \end{beamercolorbox}%
    \begin{beamercolorbox}[wd=.2\paperwidth,ht=2.5ex,dp=1ex,right]{date in head/foot}%
      \insertframenumber{} / \inserttotalframenumber\hspace*{1em}
    \end{beamercolorbox}%
  }%
  \vskip0pt%
}

\begin{document}

\begin{frame}
\titlepage
\end{frame}

\begin{frame}
\frametitle{Table of Contents (Part 1)}
\tableofcontents[sections=1-10] % Adjust the range to show the first part
\end{frame}

% Second part of the Table of Contents
\begin{frame}
\frametitle{Table of Contents (Part 2)}
\tableofcontents[sections=11-19] % Adjust the range to show the remaining sections
\end{frame}
\section{Introduction}

\begin{frame}
\centering
\Huge \textbf{Introduction}
\end{frame}

\begin{frame}{Introduction}
\begin{itemize}
\item Apricity aims to assist founders and investors by providing fast, data-driven evaluation of startup viability
\item Designed as an AI-powered software prototype that simulates expert-level startup assessment through autonomous agents
\item Employs a multi-agent architecture where each agent represents a domain expert such as Finance, Technology, Marketing, Sales, and Product
\item Integrates Large Language Models with Retrieval-Augmented Generation (RAG) to ensure informed, context-aware, and explainable predictions
\item Generates structured insights including risk analysis, strengths, weaknesses, and improvement suggestions in real time
\end{itemize}
\end{frame}


\section{Multi-Agent System (MAS)}

\begin{frame}
\centering
\Huge \textbf{Multi-Agent System (MAS)}
\end{frame}

\begin{frame}{Multi-Agent System (MAS)}
\begin{itemize}
    \item A \textbf{multi-agent system (MAS)} is a collection of AI agents working together to solve complex problems.
    \item Each agent specializes in a specific domain (e.g., market analysis, product evaluation, financial review).
    \item Agents use \textbf{Retrieval-Augmented Generation (RAG)} to access up-to-date, domain-specific information.
    \item The system is modular and \textbf{easily extendable}—new agents can be added or existing ones modified.
  
\end{itemize}
\end{frame}

\begin{frame}{Multi-Agent System (MAS)}
\begin{itemize}
    \item \textbf{Financial Analyst}: Evaluates financial viability, forecasts, cost structures, and potential ROI.
    \item \textbf{Venture Capitalist (VC)}: Assesses investment readiness, scalability, and risk factors.
    \item \textbf{CTO (Technology Officer)}: Reviews technical feasibility, innovation, and scalability of the product.
    \item \textbf{Marketing Specialist}: Analyzes target market, go-to-market strategy, and brand positioning.
    \item \textbf{Sales Manager}: Focuses on sales funnel, customer acquisition strategy, and revenue generation.
    \item \textbf{Product Analyst}: Examines product-market fit, user needs, feature set, and competitive differentiation.
\end{itemize}
\end{frame}


\begin{frame}{Multi-Agent System (MAS)}
\begin{itemize}
    \item Enables \textbf{rapid, domain-specific analysis} without requiring multiple rounds of external consultation.
    \item Provides \textbf{parallel evaluations} from multiple expert perspectives in real-time.
    \item Delivers \textbf{consistent, unbiased insights} derived from the latest data using RAG techniques.
    \item \textbf{Reduces cost and time} significantly compared to traditional manual assessments.
    \item Empowers founders to make \textbf{data-driven decisions quickly}, improving agility and market responsiveness.
\end{itemize}
\end{frame}


\section{Features}
\begin{frame}
\centering
\Huge \textbf{Features}
\end{frame}

\begin{frame}{Key Features of the System}
\begin{itemize}
    \item \textbf{AI-Powered Evaluation:} Uses large language models to assess startup viability across multiple dimensions.
    \item \textbf{Multi-Agent Architecture:} Independent agents simulate domain experts such as Financial Analyst, VC, CTO, Marketing, Sales, and Product specialists.
    \item \textbf{Retrieval-Augmented Generation (RAG):} Grounds evaluations using up-to-date external knowledge sources for improved accuracy.
    \item \textbf{Real-Time Feedback:} Generates comprehensive startup assessments within minutes.
    \item \textbf{Modular and Scalable Design:} Agents can be added, removed, or customized for different industries and evaluation criteria.
\end{itemize}
\end{frame}

\section{Purpose and Scope}

\begin{frame}
\centering
\Huge \textbf{Purpose and Scope}
\end{frame}

\begin{frame}{Purpose and Scope}
\begin{itemize}
    \item Defines the design and implementation of an AI-driven system for predicting startup success
    \item Captures both functional and non-functional requirements of the Startup Success Predictor software prototype
    \item Simulates expert-driven startup evaluation using autonomous multi-agent reasoning
    \item Focuses on software components including agent orchestration, retrieval pipeline, and evaluation logic
    \item Provides structured, explainable assessments across finance, technology, market, sales, and product domains
    \item Designed to be extensible for future integration with real-world datasets, investor tools, and decision-support platforms
\end{itemize}
\end{frame}





\section{Literature Survey}

\begin{frame}
\centering
\Huge \textbf{Literature Survey}
\end{frame}


\begin{frame}{SSFF: Investigating LLM Predictive Capabilities for Startup Succes}
\large\textbf{Xisen Wang, Yigit Ihlamur, et al.}

\normalsize
\begin{itemize}
  \item Introduces a multi-agent system combining LLMs, random forests, and neural networks.
    \item Demonstrates that founders' professional backgrounds significantly influence startup success rates.

    \item Framework includes: analyst agents, prediction block (LLM-enhanced RF model), and external market data via RAG.
    \item Addresses LLM overprediction bias using structured prompting and ensemble reasoning.
\end{itemize}

\end{frame}

\begin{frame}{Attention Is All You Need}
\large\textbf{Ashish Vaswani et al. (Google Brain, 2017)}

\normalsize
\begin{itemize}
  \item Introduced the Transformer model using only attention, removing RNNs and CNNs.
    \item Employed multi-head self-attention and positional encoding.
    \item Enabled fast, parallel training with lower computational cost.
    \item Achieved state-of-the-art results in machine translation tasks.
\end{itemize}
\end{frame}


\begin{frame}{Survey of Marketing Agents, Agent-Based Models, and Generative AI in Marketing    }
\large\textbf{Meghana Puvvadi, et al.}

\normalsize
\begin{itemize}
   \item ABMs and LLMs enable personalized, adaptive marketing beyond traditional models.
    \item ABMs simulate real-world consumer and influencer interactions.
    \item LLMs automate and personalize content across multiple channels.
    \item Hybrid ABM-LLM frameworks support scalable, real-time marketing campaigns.
\end{itemize}
\end{frame}


\begin{frame}{LLM Multi-Agent Systems: Challenges and Open Problems }
\large\textbf{Shanshan Han, et al.}

\normalsize
\begin{itemize}
 \item Identifies four agent structures: Equi-Level, Hierarchical, Nested, and Dynamic.
    \item Highlights the complexity of global planning and coordination among agents.
    \item Emphasizes the need for both local and shared memory across agents.
    \item Describes SSFF use case with shared JSON/vector memory and local agent memory.
  \item Outperforms baselines on Mind2Web across tasks and domains.
\end{itemize}
\end{frame}


\begin{frame}{Improving Startup Success with Text Analysis }
\large\textbf{Emily Gavrilenko, et al.}

\normalsize
\begin{itemize}
   \item Aims to predict startup funding outcomes using public and text-based data.
    \item Textual features (e.g., sentiment, passive voice, POS patterns) boost accuracy—especially from Twitter data.
    \item NLP on business descriptions and social media improves predictions.
    \item Combines structured (numeric) and unstructured (text) data for stronger models.
\end{itemize}
\end{frame}


\begin{frame}{Developing RAG-Based LLM Systems from PDFs: An Experience Report }
\large\textbf{Ayman Asad Khan, et al.}

\normalsize
\begin{itemize}
   \item Presents a hands-on guide for building Retrieval-Augmented Generation (RAG) systems using PDF documents as data sources.
    \item Details the full pipeline: PDF collection, text preprocessing, vector embedding, retrieval, and LLM-based response generation.
    \item Evaluates two approaches: OpenAI's Assistant API for simplicity and Llama for open-source flexibility.
    \item Addresses practical challenges such as PDF layout issues, chunking strategies, metadata handling, and error management.
\end{itemize}
\end{frame}

\begin{frame}{Survey on Evaluation of LLM-Based Agents  }
\large\textbf{Asaf Yehudai, Lilach Eden, et al.}

\normalsize
\begin{itemize}
  \item First comprehensive survey on evaluation methods for LLM-based agents.
  \item Reviews benchmarks across core skills, applications, general agents, and frameworks.
  \item Identifies gaps in cost-efficiency, robustness, and safety evaluation.
  \item Mentions emerging tools like LangSmith and Galileo for stepwise evaluation.
\end{itemize}

\end{frame}


\begin{frame}{Founder-GPT: Self-Play to Evaluate the Founder–Idea Fit}
\large\textbf{Sichao Xiong, Yigit Ihlamur}

\normalsize
\begin{itemize}
 \item Introduces a self-play multi-agent framework with Tree-of-Thought prompting for VC-style founder-idea evaluation.
    \item Simulates analyst debates refined by a critic model to derive explainable component scores.
    \item Combines embeddings of LinkedIn bios and idea text with structured academic features for similarity retrieval.
    \item Uses a composite scoring pipeline with an ethical gate—unethical ventures are automatically rejected.
\end{itemize}
\end{frame}







\begin{frame}{MaAS – Multi-agent Architecture Search via Agentic Supernet }
\large\textbf{Guibin Zhang, Luyang Niu et al.}

\normalsize
\begin{itemize}
  \item Proposes a probabilistic supernet that samples query-specific agent workflows, replacing static pipelines.
    \item Controller dynamically selects reasoning operators (e.g., I/O, CoT, Debate, ReAct) and can early-exit for simpler tasks.
    \item Achieves up to 16.9 percentage point accuracy gain and reduces inference costs by 55–94\%.
    \item Enables cost-aware, auto-designed multi-agent workflows with cross-model and cross-dataset generalization.
\end{itemize}
\end{frame}



\begin{frame}{Zero-Indexing Internet Search Augmented Generation for Large
Language Models}
\large\textbf{Guangxin He, et al.}

\normalsize
\begin{itemize}
  \item Introduces zero-indexing with dynamic search APIs instead of static corpora.
  \item Uses collaborative LLMs for search, ranking, and information extraction.
  \item Cuts token cost by up to 47 percent while improving output quality.
  \item Deployed in 01.AI’s real-time inference system.
\end{itemize}

\end{frame}

\begin{frame}{Chain of Agents: A Modular Framework for Solving Complex Tasks with LLMs}
\large\textbf{Brian Lee, Zifan Wang, et al.}

\normalsize
\begin{itemize}
   \item Proposes a modular framework where agents are composed in a chain to solve complex tasks with long-term planning.
    \item Agents operate independently, communicate via function calls, and can invoke tools or sub-agents.
    \item Demonstrates improved task performance across domains like math, coding, and decision-making.
    \item Enables interpretability and reusability by structuring agent workflows with clear modular interfaces..
\end{itemize}

\end{frame}





\begin{frame}{Literature Survey (1/2)}
\resizebox{\textwidth}{!}{%
\tiny
\begin{tabular}{|p{1.5cm}|p{1.3cm}|p{2.8cm}|p{2cm}|p{2cm}|p{2.2cm}|}
\hline
\textbf{Title} & \textbf{Authors} & \textbf{Key Findings} & \textbf{Advantages} & \textbf{Disadvantages} & \textbf{Inferences} \\
\hline

SSFF & Wang et al. (Oxford, Vela Partners) & Multi-agent system combining LLMs, random forests, neural networks for startup prediction. & Addresses LLM overprediction bias, incorporates founder background analysis. & Complexity of ensemble approach, potential overfitting to specific datasets. & Founder professional backgrounds significantly influence startup success rates \\

\hline
Attention Is All You Need & Vaswani et al. (Google Brain, 2017) & Introduced Transformer using only attention, removing RNNs/CNNs. Multi-head self-attention and positional encoding. & Fast parallel training, lower computational cost, state-of-the-art translation results, generalizes well to sequence tasks. & Not explicitly mentioned in source. & Pure attention mechanisms can replace recurrence and convolution for sequence modeling tasks. \\
\hline
Survey of Marketing Agents & Puvvadi et al. & ABMs and LLMs enable personalized, adaptive marketing. Hybrid ABM-LLM frameworks support real-time campaigns. & Personalized content across channels, real-world consumer simulation, scalable campaigns. & Privacy concerns, algorithmic bias, lack of transparency, sustainability issues. & AI-driven marketing requires balancing personalization with ethical considerations. \\
\hline
LLM Multi-Agent Systems & Han et al. & Four agent structures: Equi-Level, Hierarchical, Nested, Dynamic. Emphasis on global planning and coordination. & Local and shared memory systems, decentralized reasoning through agent debates. & Complex coordination challenges, global planning difficulties. & Multi-agent architectures need sophisticated memory and coordination mechanisms. \\
\hline
Improving Startup Success & Gavrilenko et al. (JP Morgan Chase) & Textual features from Twitter and business descriptions boost funding prediction accuracy. & Combines structured and unstructured data, NLP enhances predictions. & Limited to specific data sources, accuracy depends on text quality. & Text analysis provides valuable signals for startup investment decisions. \\
\hline


Survey on LLM Agents & Yehudai et al. (Hebrew Univ., IBM, Yale) & First comprehensive survey of LLM agent evaluation methods across four dimensions. & Systematic evaluation framework, identifies research gaps, emerging evaluation tools. & Gaps in cost-efficiency and robustness evaluation, scalability challenges. & Standardized evaluation methods needed for LLM agent development. \\
\hline

\end{tabular}}
\end{frame}

\begin{frame}{Literature Survey (2/2)}
\resizebox{\textwidth}{!}{%
\tiny
\begin{tabular}{|p{1.5cm}|p{1.3cm}|p{2.8cm}|p{2cm}|p{2cm}|p{2.2cm}|}
\hline
\textbf{Title} & \textbf{Authors} & \textbf{Key Findings} & \textbf{Advantages} & \textbf{Disadvantages} & \textbf{Inferences} \\




\hline
Founder-GPT & Xiong, Ihlamur (Oxford, Vela Partners) & Self-play multi-agent framework with Tree-of-Thought for founder-idea evaluation. & Explainable scoring, ethical gate, interpretable evaluation using only LinkedIn and idea input. & Limited to specific input types, depends on LinkedIn data quality. & Multi-agent self-play enables sophisticated founder-idea fit assessment. \\
\hline
MaAS & Zhang et al. (ICML 2025) & Probabilistic supernet samples query-specific agent workflows dynamically. & 16.9pp accuracy gain, 55-94\% cost reduction, cross-model generalization. & Complex architecture, requires training supernet controller. & Dynamic workflow selection optimizes both performance and efficiency. \\

\hline
Zero-Indexing Internet Search & He et al. (HKUST, Fudan, 01.AI, Tsinghua) & Uses zero-indexing paradigm leveraging real-time web search APIs with collaborative LLMs for query parsing, re-ranking, and content extraction. & Achieves higher quality content, reduces input tokens by 21–47\%, and is production-deployed for real-time up-to-date generation. & Dependent on external search APIs; real-time search latency and noise in web data pose challenges. & Enables dynamic, fresh online knowledge integration without maintaining static indexes, improving retrieval relevance and efficiency. \\

\hline


Chain of Agents & Lee et al. (Stanford University) & Modular framework where agents are composed in chains to solve complex, multi-step tasks. & Improves task performance in math, coding, and planning; promotes modularity and interpretability. & Requires careful design of agent interfaces; potential communication overhead. & Modular agent chaining enables scalable and reusable LLM-driven systems for complex reasoning. \\



\hline
Developing RAG-Based LLM Systems from PDFs & Khan et al. (Tampere University) & Step-by-step guide to building RAG systems using PDFs; compares OpenAI API and Llama approaches. & Real-time retrieval, transparent responses, supports both proprietary and open-source tools. & Complex PDF layouts pose extraction challenges; OpenAI vector store raises privacy concerns. & RAG improves factual accuracy and usability in knowledge-intensive domains. \\
\hline
\end{tabular}}
\end{frame}

%------------------------------------------------
%------------------------------------------------
\section{Problem Statement}

\begin{frame}
\centering
\Huge \textbf{Problem Statement}
\end{frame}

\begin{frame}{Problem Statement}
\begin{itemize}
    \item Startup evaluation today relies heavily on manual expert reviews.
    \item These processes are slow, expensive, and difficult to scale.
    \item Early-stage founders often lack access to structured, expert-level feedback.
\end{itemize}
\end{frame}

\section{Proposed Solution}

\begin{frame}
\centering
\Huge \textbf{Proposed Solution}
\end{frame}

\begin{frame}{Proposed Solution}
\begin{itemize}
    \item An AI-driven multi-agent system that simulates expert startup evaluation.
    \item Each agent represents a domain expert (finance, product, market, etc.).
    \item Produces structured insights instead of a single opaque prediction.
\end{itemize}
\end{frame}

\section{Project Progress Overview}

\begin{frame}
\centering
\Huge \textbf{Project Progress Overview}
\end{frame}

\begin{frame}{Project Progress Overview (Current)}
\begin{itemize}
    \item Overall Status: \textbf{$\sim$90\%} core MVP pipeline implemented
    \item Frontend: Guided startup workspace (templates + rich-text editor), autosave, PDF upload + viewer
    \item Backend: FastAPI endpoint runs \textbf{multi-agent analysis} with RAG chunking + retrieval
    \item Web enrichment: \textbf{Exa API} competitor discovery implemented (structured competitor list)
    \item Report: Results rendered in UI + \textbf{PDF export implemented} via \textbf{html2pdf.js}
\end{itemize}
\end{frame}

\section{System Flow}

\begin{frame}
\centering
\Huge \textbf{System Flow Diagram}
\end{frame}

\begin{frame}{End-to-End System Flow}
\centering
\IfFileExists{Block.png}{
\includegraphics[
    width=0.9\textwidth,
    height=0.75\textheight,
    keepaspectratio
]{Block.png}

\vspace{0.5em}
\small Flow diagram showing end-to-end data movement across system modules.
}{
\fbox{\parbox{0.85\textwidth}{Flow diagram placeholder: \texttt{Block.png}}}
}
\end{frame}

\begin{frame}{System Flow Explanation (1/2)}
\small
\begin{enumerate}
    \item \textbf{User Input -> Workspace (Frontend):}
    User fills guided startup sections using artifact templates (rich-text editor) and can upload a PDF for reference.

    \item \textbf{Workspace -> Local Persistence:}
    Documents, files, and view state are autosaved to \textbf{IndexedDB} (debounced) for reliability across refreshes.

    \item \textbf{Workspace -> "Develop a View":}
    Frontend aggregates all document content into a single analysis prompt and creates a new \texttt{viewId} with \texttt{pending} status.

    \item \textbf{View Page -> Analysis Request:}
    UI shows progress stages and sends the aggregated prompt to the internal API endpoint \texttt{/api/views}.

    \item \textbf{API Route -> Backend Forwarding Setup:}
    Next.js API validates the request and forwards it to the Python backend endpoint \texttt{/view}.
\end{enumerate}
\end{frame}

\begin{frame}{System Flow Explanation (2/2)}
\small
\setcounter{enumi}{5}
\begin{enumerate}
    \item \textbf{Next.js API -> FastAPI Backend:}
    The prompt is posted to \texttt{POST /view} for analysis.

    \item \textbf{Backend -> RAG Context Preparation:}
    Input is chunked, embedded (OpenAI embeddings), and indexed in \textbf{Chroma} to enable retrieval during agent reasoning.

    \item \textbf{Backend -> Multi-Agent Evaluation:}
    \textbf{LangGraph} runs domain agents (Finance, VC, CTO, Marketing, Product) in parallel using the shared retriever.

    \item \textbf{Backend -> Web Enrichment (Competitors):}
    \textbf{Exa API} is used to discover and enrich top competitors into a structured competitor dataset (disabled if \texttt{EXA\_API\_KEY} is missing).

    \item \textbf{Backend -> Structured JSON Response:}
    Backend returns agent analyses, competitor search status, and competitor list as JSON.

    \item \textbf{Frontend -> Results Rendering \& Storage:}
    The view is marked \texttt{completed} (or \texttt{error}), results are rendered section-wise, and stored in the workspace state for later access.

    \item \textbf{Frontend -> Report Export:}
    User exports a styled A4 PDF report using \textbf{html2pdf.js} (markdown rendering supported for agent outputs).
\end{enumerate}
\end{frame}

%------------------------------------------------
\section{Novelty of the Project}
%------------------------------------------------

\begin{frame}
\frametitle{Novelty of Project}
    \textbf{Novelty}
    \begin{itemize}

\item Uses a multi agent system where each agent focuses on a different aspect of the startup, combining their insights into one coherent analysis.
\item Connects to live external data sources to include current market and industry information in every evaluation.
\item Replaces traditional numeric scoring with qualitative, human readable predictions that explain reasoning clearly and build trust in the decision process.
\item Emphasizes transparent, language based reasoning instead of complex numeric scores for better interpretability.
\item Adopts a modular design that can evolve by adding new agents or data inputs as startup evaluation methods advance.
    \end{itemize}
\end{frame}
%------------------------------------------------
\section{Objectives}
%------------------------------------------------

\begin{frame}
\frametitle{Project Objectives}
    \textbf{Objectives}
    \begin{itemize}
        \item Develop an intelligent multi-agent system to evaluate startup ideas comprehensively
        \item Enable automated analysis across financial, technical, marketing, and product domains
        \item Provide data-driven success predictions with actionable recommendations
        \item Reduce evaluation time and costs compared to traditional expert consultations
        \item Deliver a user-friendly web interface for startup founders and investors to assess business viability
    \end{itemize}
\end{frame}

%------------------------------------------------
\section{Software Requirements Specification}
\begin{frame}
\begin{center}
    \huge \textbf{Software Requirements Specification}
\end{center}
\end{frame}
%------------------------------------------------




% Product Overview
\begin{frame}{Product Overview}
    \textbf{Purpose:} To provide a functional description of a system that evaluates startup ideas using AI agents.\\[6pt]
    \textbf{Scope:}
    \begin{itemize}
        \item Analyzes user-submitted startup ideas across financial, technical, and marketing domains.
        \item Predicts the potential for startup success.
        \item Generates detailed reports with insights and recommendations.
    \end{itemize}
\end{frame}

% Multi-Agent System
\begin{frame}{Multi-Agent System (MAS)}
    The system consists of multiple specialized AI agents working collaboratively:
    \begin{itemize}
        \item \textbf{Financial Analyst:} Evaluates financial viability and ROI.
        \item \textbf{CTO:} Reviews technical feasibility and innovation.
        \item \textbf{Marketing Specialist:} Analyzes target market and branding.
        \item \textbf{Product Analyst:} Assesses market fit and competitive edge.
    \end{itemize}
    These agents combine their analyses to produce a unified prediction report.
\end{frame}

% Overall Description
\begin{frame}{Overall Description}
    \textbf{Product Perspective:}
    \begin{itemize}
        \item Replaces traditional expert consultations with a fast, automated system.
        \item Uses a \textbf{multi-agent AI architecture} for domain-specific evaluation.
    \end{itemize}
    \vspace{6pt}
    \textbf{Product Functions:}
    \begin{itemize}
        \item Data ingestion and augmentation using RAG.
        \item Parallel analysis by expert agents.
        \item Consolidated report generation.
    \end{itemize}
\end{frame}

% User Classes
\begin{frame}{User Classes and Characteristics}
    \begin{itemize}
        \item \textbf{Startup Founders:} Validate ideas before funding.
        \item \textbf{Product Managers:} Assess market fit and risks.
        \item \textbf{Developers:} Test technical feasibility.
        \item \textbf{Investors:} Evaluate early-stage ideas for potential success.
    \end{itemize}
\end{frame}

% Operating Environment
\begin{frame}{Operating Environment}
    \begin{itemize}
        \item Web-based application accessible via browsers (Chrome, Firefox, Safari).
        \item Backend built with \textbf{Python + LangChain}.
        \item Frontend built with \textbf{React}.
        \item Requires API connectivity with external LLM services.
    \end{itemize}
\end{frame}

\subsection{External Interface Requirements}
% External Interface Requirements
\begin{frame}{External Interface Requirements}
    \textbf{User Interfaces:}
    \begin{itemize}
        \item Intuitive web UI for data input and report visualization.
    \end{itemize}
    \vspace{4pt}
    \textbf{Hardware Interfaces:}
    \begin{itemize}
        \item Standard PC or mobile device with web browser.
    \end{itemize}
    \vspace{4pt}
    \textbf{Software Interfaces:}
    \begin{itemize}
        \item Python backend with LLM APIs.
        \item Database for storing reports and metadata.
    \end{itemize}
    \vspace{4pt}
    \textbf{Communication:} Secure HTTPS protocol.
\end{frame}

% Hardware and Software Requirements
\subsection{Hardware and Software Requirements}
\begin{frame}{Hardware and Software Requirements}
    \textbf{Hardware:}
    \begin{itemize}
        \item Server with 8GB+ RAM and dual-core CPU.
        \item Client: Any modern computer or smartphone.
    \end{itemize}
    \vspace{6pt}
    \textbf{Software:}
    \begin{itemize}
        \item Python 3.8+ for backend.
        \item React for frontend.
        \item Database for user data.
        \item LLM API key for agent reasoning.
    \end{itemize}
\end{frame}

\subsection{Functional Requirements}
\begin{frame}{Functional Requirements}
    \begin{enumerate}
        \item \textbf{User Input \& Data Handling:} Collects startup details and fetches additional context via RAG.
        \item \textbf{Multi-Agent Analysis:} Domain-specific agents evaluate the idea collaboratively.
        \item \textbf{Report Generation:} Produces detailed success prediction and recommendations.
        \item \textbf{User Interaction:} Provides progress updates and notifications.
        \item \textbf{Security \& Privacy:} Encrypts data and ensures confidentiality.
    \end{enumerate}
\end{frame}

\subsection{Nonfunctional Requirements}
\begin{frame}{Nonfunctional Requirements}
    \textbf{Performance:} Fast and accurate analysis delivery.\\[6pt]
    \textbf{Safety:} Prevents misinformation by validating sources.\\[6pt]
    \textbf{Security:} All data encrypted; no permanent storage of confidential info.\\[6pt]
    \textbf{Quality Attributes:}
    \begin{itemize}
        \item Modularity for agent addition/removal.
        \item Reliability and transparency in predictions.
    \end{itemize}
\end{frame}

\section{Software Design Document}
%------------------------------------------------
\begin{frame}
\begin{center}
    \huge \textbf{Software Design Document}
\end{center}
\end{frame}



%------------------------------------------------
\subsection{Frontend}
\begin{frame}{Frontend}
    The frontend serves as the user interface, enabling data input and interaction with the system.\\[8pt]
    \textbf{Key Functions:}
    \begin{itemize}
        \item \textbf{Input Form:} Collects essential business data — idea, founders, market, strategy, sales, and finance.
        \item \textbf{Evaluation Trigger:} Initiates the multi-agent analysis and tracks progress in real time.
        \item \textbf{Result Dashboard:} Displays consolidated agent results and overall evaluation insights.
        \item \textbf{Report Generation:} Summarizes agent perspectives, risks, suggestions, and evaluation scores.
        \item \textbf{User Interaction:} Provides a smooth and intuitive interface for easy data submission and feedback.
    \end{itemize}
\end{frame}
%------------------------------------------------

%------------------------------------------------
\subsection{Backend}
\begin{frame}{Backend}
    The backend powers the analytical intelligence and data processing components of the system.\\[8pt]
    \textbf{Core Components:}
    \begin{itemize}
        \item \textbf{API Orchestration:} Manages communication between frontend, agents, and storage.
        \item \textbf{Knowledge Retrieval:} Gathers and summarizes relevant market and product data from the web.
        \item \textbf{Vector Repository:} Stores semantic embeddings for context-aware retrieval and matching.
        \item \textbf{Multi-Agent Coordination:} Synchronizes agents handling market, product, and founder evaluations.
        \item \textbf{Prediction Engine:} Generates interpretable scores and insights from structured and unstructured data.
        \item \textbf{Data Management:} Handles persistent storage, versioning, and retrieval for reproducible analyses.
    \end{itemize}
\end{frame}

%------------------------------------------------
\subsection{Input Data}
\begin{frame}{Input Data}
    \begin{itemize}
        \item \textbf{Business Idea} – Core concept or innovation behind the project.
        \item \textbf{Founder Information} – Details about the team, expertise, and roles.
        \item \textbf{Market Study} – Research data, trends, and competitor analysis.
        \item \textbf{Marketing Strategy} – Target audience, positioning, and promotion plans.
        \item \textbf{Sales Information} – Sales approach, channels, and customer acquisition plan.
        \item \textbf{Financial Information} – Funding, expenses, and projected revenue details.
    \end{itemize}
\end{frame}
%------------------------------------------------
%------------------------------------------------

%------------------------------------------------
\subsection{Use Case Diagram}
%------------------------------------------------
\begin{frame}
\frametitle{Use Case Diagram}
\centering
\includegraphics[width=0.75\linewidth]{usecase2.drawio.png} % replace with actual file
\end{frame}




\section{Data Flow Diagram}
%------------------------------------------------
\begin{frame}
\begin{center}
    \huge \textbf{Data Flow Diagram}
\end{center}
\end{frame}
%------------------------------------------------
\subsection{Level 0}
%------------------------------------------------
\begin{frame}
\frametitle{Data Flow Diagram}
\textbf{LEVEL 0:}
\begin{figure}
    \includegraphics[width=\linewidth]{Level0.png}
    \caption{Level 0 Data Flow Diagram}
\end{figure}
\end{frame}


%------------------------------------------------
\subsection{Level 1}
%------------------------------------------------
\begin{frame}
\frametitle{Data Flow Diagram}
\textbf{LEVEL 1:}
\begin{figure}
    \includegraphics[width=1\linewidth]{dfd1.jpeg}
    \caption{Level 1 Data Flow Diagram}
\end{figure}
\end{frame}

%------------------------------------------------
\subsection{Level 2}
%------------------------------------------------
\begin{frame}
\frametitle{Data Flow Diagram}
\textbf{LEVEL 1.1:}
\begin{figure}
    \includegraphics[width=1\linewidth]{dfdPics/Untitled Diagram.drawio(6).png}
    \caption{Level 1.1 Data Flow Diagram}
\end{figure}
\end{frame}

\begin{frame}
\frametitle{Data Flow Diagram}
\textbf{LEVEL 1.2:}
\begin{figure}
    \includegraphics[width=0.85\linewidth]{dfd1_2.png}
    \caption{Level 1.2 Data Flow Diagram}
\end{figure}
\end{frame}


\begin{frame}
\frametitle{Data Flow Diagram}
\textbf{LEVEL 1.3:}
\vspace{10}

\begin{figure}
    \includegraphics[width=1\linewidth]{level1_3.jpeg}
    \caption{Level 1.3 Data Flow Diagram}
\end{figure}
\end{frame}





%<-- Darwin was here 




\section{System Architecture}
%------------------------------------------------
\begin{frame}
\begin{center}
    \huge \textbf{System Architecture}
\end{center}
\end{frame}
%------------------------------------------------
\begin{frame}
\frametitle{System Architecture}
\textbf{Block Diagram}
\begin{figure}
    \includegraphics[width=0.5\linewidth]{flowchart final (1).png}
    \caption{Block Diagram}
\end{figure}
\end{frame}


%------------------------------------------------
\section{Module Description}
%------------------------------------------------

\begin{frame}{Input Module}
\small
\begin{itemize}
    \item \textbf{Purpose / Functions:} Create and edit startup artifacts, upload PDFs, autosave workspace state, build analysis payload.
    \item \textbf{Input / Output:} Artifact text + PDFs -> saved workspace + consolidated startup context.
    \item \textbf{Libraries:} React / Next.js, TipTap, idb-keyval.
\end{itemize}
\end{frame}

\begin{frame}{Web Searcher Module}
\small
\begin{itemize}
    \item \textbf{Purpose / Functions:} Generate search sentence, fetch competitor data, normalize response.
    \item \textbf{Input / Output:} Startup idea text -> competitor list + search status.
    \item \textbf{Libraries:} httpx, ChatOpenAI, ChatPromptTemplate, StrOutputParser.
\end{itemize}
\end{frame}

\begin{frame}{Agent Module}
\small
\begin{itemize}
    \item \textbf{Purpose / Functions:} Chunk context, build retriever, run 5 agents, aggregate final result.
    \item \textbf{Input / Output:} Combined startup context -> agent analyses + overall evaluation.
    \item \textbf{Libraries:} langgraph, ChatOpenAI, OpenAIEmbeddings, Chroma, RecursiveCharacterTextSplitter.
\end{itemize}
\end{frame}

\begin{frame}{Report Generation Module}
\small
\begin{itemize}
    \item \textbf{Purpose / Functions:} Render results in UI and export final PDF.
    \item \textbf{Input / Output:} Analysis JSON -> report view + PDF export.
    \item \textbf{Libraries:} React / Next.js, html2pdf.js, react-markdown.
\end{itemize}
\end{frame}

\section{Progress Summary}

\begin{frame}{Conclusion \& Next Steps}
\begin{itemize}
    \item Project is on track with a \textbf{modular, incremental implementation} toward an MVP
    \item Core modules are \textbf{implemented and working} end-to-end:
    \begin{itemize}
        \item Input workspace (\textbf{Next.js + React}): guided templates, rich-text editor, PDF upload/viewer, autosave
        \item Web enrichment (\textbf{Exa API}): competitor discovery + structured competitor dataset
        \item Multi-agent evaluation (\textbf{LangGraph + OpenAI}): domain agents with RAG-based context retrieval
        \item Report output: results UI + \textbf{PDF export} implemented (html2pdf.js)
    \end{itemize}
    \item Current focus: \textbf{performance and quality optimization} (prompt tuning, token efficiency, enrichment reliability, UI/UX polish)
    \item Next milestone: improved citations/source attribution + caching/rate limits + more polished, consistent PDF reports
\end{itemize}
\end{frame}

\section{Conclusion}


\begin{frame}
\frametitle{Conclusion}
\begin{itemize}
   \item The uncertainty of startup evaluation is transformed into clarity and direction.
\item Guesswork and scattered opinions give way to structured understanding and insight.
\item The platform guides founders to view their ideas through the power of experience and data.
\item  It shapes the future with confidence and precision instead of mere prediction.
\item  The vision of giving every founder expert level understanding becomes a reality.


\end{itemize}
\end{frame}

\section{References}

\begin{frame}
\frametitle{References}
\footnotesize{

[1] A. Vaswani, N. Shazeer, N. Parmar, et al., “Attention Is All You Need,” in \textit{Advances in Neural Information Processing Systems}, 2017. \\

[2] M. Puvvadi, S. K. Arava, A. Santoria, et al., “Survey of Marketing Agents, Agent-Based Models, and Generative AI in Marketing,” \textit{arXiv preprint} arXiv:2401.01234, 2024. \\

[3] S. Han, Q. Zhang, Y. Yao, W. Jin, and Z. Xu, “LLM Multi-Agent Systems: Challenges and Open Problems,” \textit{arXiv preprint} arXiv:2403.04567, 2024. \\

[4] E. Gavrilenko, F. Khosmood, M. Rastad, et al., “Improving Startup Success with Text Analysis,” \textit{JP Morgan Chase Technical Report}, 2023. \\

[5] A. Khan and M. V. Mäntylä, “Developing RAG-Based LLM Systems from PDFs,” \textit{Tampere University}, 2024. \\
}
\end{frame}

\begin{frame}
\frametitle{References}
\footnotesize{
[6] A. Yehudai, L. Eden, A. Li, et al., “Survey on Evaluation of LLM-Based Agents,” \textit{arXiv preprint} arXiv:2402.07890, 2024. \\

[7] S. Xiong and Y. Ihlamur, “Founder-GPT: Self-Play to Evaluate the Founder–Idea Fit,” \textit{arXiv preprint} arXiv:2404.01111, 2024. \\

[8] G. Zhang, L. Niu, et al., “MaAS – Multi-Agent Architecture Search via Agentic Supernet,” in \textit{Proceedings of ICML}, 2025. \\

[9] P. Yin, J. Liu, et al., “Zero-Indexing: Decoupling Planning and Execution in LLMs,” Stanford, UC Berkeley, Microsoft Research, 2024. \\

[10] K. Lee, J. Chen, et al., “Chain of Agents: Modular Composition for Complex Reasoning,” \textit{Stanford University}, 2024. \\

[11] X. Wang, Y. Ihlamur, and F. Alican, “SSFF: Investigating LLM Predictive Capabilities for Startup Success,” \textit{arXiv preprint} arXiv:2405.02777, 2024. \\

}
\end{frame}



\begin{frame}
\centering
\Huge \textbf{Thank You}
\end{frame}
\end{document}
